\title{LongRoPE: Extending LLM Context Window Beyond 2 Million Tokens}

\begin{abstract}
Expanding the context window in large language models (LLMs) is highly advantageous. Nonetheless, challenges such as the expensive nature of fine-tuning, limited availability of lengthy textual data, and the emergence of disruptive values from unfamiliar token positions constrain existing context window extensions to approximately 128k tokens.
\end{abstract}

This work presents LongRoPE, a method that, for the first time, expands the context window of pre-trained LLMs up to an unprecedented \textbf{2048k} tokens, requiring no more than 1k fine-tuning steps at training lengths up to 256k, all while preserving their performance at the default short context window. This advancement is made possible by three major contributions: (i) we discover and utilize two distinct types of positional interpolation non-uniformities through an optimized search, yielding superior fine-tuning initialization and permitting an 8$\times$ extension even without further fine-tuning; (ii) a novel progressive extension approach is adopted, wherein a model is initially fine-tuned to handle 256k tokens before a subsequent positional interpolation enlarges its context window to 2048k; (iii) LongRoPE is subsequently recalibrated at an 8k context length to restore its performance on shorter contexts. Comprehensive empirical studies using LLaMA2 and Mistral on multiple benchmarks verify the robustness of our technique. Models extended by LongRoPE largely retain their original architectures, with only minor updates to the positional embedding scheme, and are compatible with most existing optimization strategies. The implementation will be released at \texttt{https://github.com/microsoft/LongRoPE}

\section{Introduction}

Although Large Language Models (LLMs) have demonstrated impressive capabilities across a range of tasks~\cite{gpt4,llama2}, they are constrained by relatively small context window sizes—for example, LLaMA2 is restricted to 4096 tokens~\cite{llama2}. When inputs exceed this window, LLMs experience decreased effectiveness since they encounter sequence lengths outside their training distribution. This limitation becomes particularly problematic in critical use cases such as in-context learning involving many examples~\cite{Huang2023FewerIM} and for LLM agents~\cite{park2023generative,madaan2023selfrefine}.

Recent studies have demonstrated that the context window of a pre-trained LLM can be scaled up to approximately 128k tokens by employing fine-tuning techniques on extended text sequences~\cite{longlora,pi,yarn,zhang2024soaring,liu2023scaling}. Nevertheless, three primary challenges remain when attempting to lengthen the context window further. First, the presence of previously unseen position indices that have not been trained on can result in a proliferation of problematic values, causing distribution shifts and hindering the convergence of fine-tuning~\cite{pi}. This obstacle is especially pronounced when expanding the context from 4k to over 1000k, as more than 90\% of position indices are then newly introduced. Second, effective fine-tuning typically necessitates text samples matching the target length, yet current data repositories contain few examples surpassing 1000k in length. Additionally, the process of training on such extensive texts incurs substantial computational costs, demanding significant time and GPU usage. Third, as the context window grows extremely large, attention mechanisms become less focused, distributing attention across a vast array of token positions; this dilutes attention density and ultimately reduces model performance on tasks with shorter input sequences~\cite{pi}.

A commonly adopted strategy to address the initial challenge involves interpolating RoPE positional embeddings~\cite{rope,pi}, wherein the new positional indices are rescaled into the interval used during pretraining (see Fig.\ref{fig:overview}). Position Interpolation (PI)~\cite{pi} performs a linear interpolation of the rotary angles in RoPE according to the scaling ratio. The NTK method~\cite{ntk,dynamicntk} proposes an approach where interpolation and extrapolation are applied differently across various dimensions of RoPE. YaRN~\cite{yarn} further divides the RoPE dimensions into three distinct frequency-based categories, utilizing extrapolation, NTK, and linear interpolation methods for each, respectively. Nonetheless, the positional embedding within the Transformer framework is characterized by an inherently \emph{complex and non-uniform information entropy}. Existing solutions do not take full advantage of this subtle variability, which results in information degradation and imposes constraints on the achievable context window size.

Section~\ref{sec:analysis} presents empirical evidence supporting two primary insights: \textbf{(1)} Optimal positional interpolation must account for non-uniformities stemming from both the diverse RoPE dimensions and the distribution of token positions. Notably, smaller RoPE dimensions as well as earlier token indices necessitate less interpolation, but what constitutes the optimal method is contingent on the desired length extension. \textbf{(2)} Incorporating these non-uniform characteristics into positional interpolation enables preservation of critical information intrinsic to the original RoPE—especially within salient dimensions and specific token positions. This approach mitigates information loss incurred by the interpolation process and yields more favorable initial conditions for subsequent fine-tuning. Additionally, it supports achieving up to an 8$\times$ extension of context length without requiring fine-tuning.

Prompted by these observations, we propose \textbf{LongRoPE}, a robust approach for extending the LLM context window to accommodate over 2 \emph{million} tokens. LongRoPE centers on three principal innovations. To begin with, {LongRoPE} takes full advantage of multidimensional non-uniformities in positional interpolation. Specifically, it determines optimal rescaling factors for RoPE's rotational parameters across each dimension, relative to token positions. Since the space of possible rescaling factors grows exponentially with the degree of context window extension, {LongRoPE} employs an evolutionary search algorithm enhanced by two optimization strategies to improve the search process's efficiency. Figure~\ref{fig:overview} presents an illustrative example of a rescaled RoPE obtained via this method.

Subsequently, {LongRoPE} adopts a resource-efficient, incremental extension approach, enabling the model to support a 2048k context window without directly fine-tuning on ultra-long texts, which are rare and difficult to acquire. Initially, the method involves probing a context length of 256k on the existing pre-trained LLM, followed by fine-tuning with this window size. Leveraging our non-uniform positional interpolation—which permits an 8$\times$ expansion without further fine-tuning—we then perform a secondary parameter search to identify appropriate RoPE rescaling factors using the previously fine-tuned model. This pipeline successfully extends the context window to 2048k tokens for both LLaMA2 and Mistral~\cite{mistral}.

In order to address potential performance loss on the original, shorter context windows, {LongRoPE} maintains the adaptation of RoPE rescale parameters within the extended LLM. Drawing parallels to the process of expanding the context window from 256k to 2048k, our approach also involves scaling down to 4k and 8k context windows on the LLM fine-tuned at 256k, employing our search algorithm to limit positional interpolation. When conducting inference with sequence lengths below 8k, we apply the rescale factors identified during this search to update RoPE accordingly.

A comprehensive set of experiments conducted on multiple LLMs and a range of long-context tasks validate the superiority of our approach. Our findings indicate that LongRoPE consistently preserves low perplexity across evaluation lengths spanning from 4k up to 2048k, attains more than 90\% accuracy in passkey retrieval, and maintains performance that matches state-of-the-art results on established benchmarks built for a 4096-token context window. Moreover, LongRoPE is compatible with any LLM utilizing RoPE embeddings. We plan to make both our implementation and the LongRoPE-2048k model variants publicly available.

\section{Non-uniformity in Positional Interpolation}
\label{sec:analysis}

\subsection{Preliminary}

Transformer architectures depend on explicit positional cues—commonly realized through position embeddings—to encode the sequential arrangement of input tokens. In this study, we concentrate on the RoPE~\cite{rope} position embedding technique, a prevalent choice among contemporary LLMs. The RoPE encoding for a token situated at position $n$ can be succinctly expressed as:

\begin{equation}
		\fontsize{7.8}{7.8} \selectfont
	\label{eq:rope}
	\left[ cos(n\theta_0), sin(n\theta_0), cos(n\theta_1), \cdots, cos(n\theta_{d/2-1}), sin(n\theta_{d/2-1}) \right]   
\end{equation}

Here, $d$ denotes the dimensionality of the embedding, $n\theta_i$ refers to the rotary positional angle of a token at index $n$, and $\theta_i = \theta^{-2i/d}$ specifies the frequencies for rotation. Typically, the value assigned to $\theta$ in RoPE is set to 10000.

\textbf{\emph{Context window extension ratio $s$} and positional interpolation}. We define $s$ as  the ratio of extended context length $L'$ to the original length $L$: $s=\frac{L'}{L}$. 

In order to increase the context window from $L$ to $L'$, existing positional interpolation techniques typically involve reducing the rotational frequencies $\theta_i$ in accordance with the extension ratio $s$. By defining $\beta = \theta^{2/d}$ and letting $\lambda$ be the corresponding rescaling factor linked to $s$, we can express these various positional interpolation strategies under a single formulation:

\begin{equation}	
	\fontsize{7.5}{7.5} \selectfont
	\label{eq:unified_rope}
	\begin{aligned}
		\left[cos\left(\frac{n}{\lambda(\beta)^0}\right), sin \left(\frac{n}{\lambda(\beta)^0}\right), cos\left(\frac{n}{\lambda(\beta)^1}\right), \cdots,  sin \left(\frac{n}{\lambda(\beta)^{d/2-1}}\right) \right]   
	\end{aligned}
\end{equation}

\textbf{Linear positional interpolation (PI)}. PI~\cite{pi} employs a strategy of linearly interpolating position indices when generating sequences beyond the model’s original training length. Given a desired extension ratio $s$, all rotational angles associated with positional encodings are uniformly scaled by $\lambda = s$ in every RoPE dimension. This uniform compression condenses positional representations, thereby reducing the model's ability to discriminate among tokens that appear close to one another in the input. As a result, PI’s performance generally diminishes at larger extension ratios.

\textbf{NTK-based interpolation and extrapolation}. ~\cite{ntk,dynamicntk} reinterpret RoPE within the framework of information representation, leveraging Neural Tangent Kernel (NTK) theory~\cite{jacot2018neural,tancik2020fourier}. To address the issue of positional crowding present in PI, they advocate for a redistribution of the interpolation effects among the various dimensions of RoPE. Their approach involves reducing the scaling applied to the lower (higher-frequency) dimensions while increasing it for the higher (lower-frequency) ones, thereby facilitating both interpolation and extrapolation over positions, where $\lambda=s^{i}$. The more advanced dynamic NTK~\cite{dynamicntk} further tunes the extension factor dynamically at each position according to the input sequence length. In contrast to PI—which requires model fine-tuning—NTK-based strategies make it possible to expand the context window even without fine-tuning, though this typically enables at most a 4$\times$ extension.

\textbf{YaRN}~\cite{yarn} brings notable advances in positional interpolation by partitioning the RoPE dimensions according to their frequency content into three separate categories, each with a specialized interpolation method. For high frequency components, extrapolation is applied ($\lambda$=1); low frequency components are interpolated linearly (PI); and intermediate frequencies are treated using NTK. The primary innovation of YaRN is its frequency-based dimension grouping in RoPE. However, the selection of these groups currently relies on empirical heuristics determined by humans, which may not generalize well and could lead to less-than-ideal outcomes for emerging LLM architectures.

\subsection{Study on Non-uniform Positional Interpolation}

Drawing upon the approaches of NTK and YaRN, we observe that their improvements stem largely from exploiting non-linearity, particularly through treating RoPE dimensions with varying frequencies differently to better tailor both interpolation and extrapolation. Despite these advancements, existing non-linearities predominantly depend on hand-crafted design choices. This prompts two central questions: (1) Does the prevailing method for positional interpolation represent the best achievable approach? (2) Might there exist other, yet untapped, forms of non-linearity?

In addressing these questions, we employ evolution search (refer to Sec.~\ref{sec:method}) to identify improved non-uniform positional interpolation strategies for LLaMA2-7B. The search process is directed by perplexity metrics, evaluated on 5 randomly selected examples from the PG19~\cite{pg19} validation dataset. Our experimental investigations yield several important insights.

\textbf{Finding 1}: \textit{RoPE dimensions exhibit substantial non-uniformities, which are not effectively handled by current positional interpolation methods.}

We conduct a search to identify the optimal $\lambda$ value for each RoPE dimension as specified in Eq.~\ref{eq:unified_rope}. Table~\ref{tbl:exp1} presents a comparison of perplexity scores for LLaMA2-7B across various methods on the PG19 and Proof-pile~\cite{proof-pile} test datasets, all evaluated without any fine-tuning. The solution identified through our search process yields notable performance gains, indicating that existing approaches based on linear (PI) and non-uniform (Dynamic-NTK and YaRN) interpolation do not achieve optimal results. In particular, YaRN performs worse than PI and NTK on PG19, which can be attributed to its inability to realize the full target context window in language models that have not undergone fine-tuning. For instance, the perplexity for YaRN rises sharply after a 7k token input in an 8k context setting.

In our investigation, the rescaling coefficients $\lambda$ in Eq.~\ref{eq:unified_rope} are not uniform, contrasting with the constant scaling $s$ applied in the methods PI, the calculation used in NTK, and the group-level approach in YaRN. The introduction of these varied scaling factors leads to notable enhancements in language modeling performance—specifically, lower perplexity—for LLaMA2 across extended context windows of 8k and 16k, all achieved without the need for additional fine-tuning. This improvement arises because the constructed positional embedding better maintains the characteristics of the original RoPE, especially along the critical dimensions, making it easier for the LLM to differentiate between tokens that are positioned closely together.

\textbf{Finding 2}: \textit{RoPE for the initial tokens in the input sequence should be extrapolated with less interpolation.} 

For the first $\hat{n}$ tokens in input sequences, we propose that reducing the degree of interpolation performed by RoPE is beneficial. These tokens typically receive higher attention scores and thus hold greater significance within attention layers, an effect documented in studies such as Streaming LLM~\cite{streamingllm} and LM-Infinite~\cite{han2023lminfinite}. To test this, we expand the context window to 8k and 16k using PI and NTK, while excluding the initial $\hat n$ tokens (where $\hat n$ = 0, 2, ..., 256) from interpolation. Setting $\hat n$ to 0 is equivalent to the baseline PI and NTK methods. As shown in Table~\ref{tbl:tokenposition}, our findings are twofold: \textbf{(1)} maintaining the original positional information for the initial tokens (i.e., skipping interpolation for these tokens) consistently enhances model performance; and \textbf{(2)} the ideal value of $\hat n$ varies based on the desired context length extension.

\textbf{Finding 3}: \textit{Non-uniform positional interpolation effectively extends LLM context window in both fine-tuning and non-fine-tuning settings.} 

Our experiments demonstrate that the use of our optimized non-uniform position interpolation yields substantial gains in extension capabilities at 8k and 16k context lengths, even without additional fine-tuning. However, to support even greater context lengths, fine-tuning becomes necessary. Accordingly, we fine-tune LLaMA2-7B with our custom-designed RoPE configuration for a 64k context window (details provided in the Appendix). As displayed in Table~\ref{tbl:finetune}, our approach markedly surpasses both PI and YaRN methods, not only in the pre-fine-tuned state but also following fine-tuning of LLaMA2-7B. These superior results can be attributed to our strategic use of non-uniform positional interpolation, which minimizes information degradation and offers superior initialization for the fine-tuning process.

\textbf{Summary}. In this work, we identify two types of non-uniformity—in the dimensions of RoPE and in token positions. By capitalizing on these aspects within positional interpolation, we achieve notable improvements in large language model context window extension.

\section{LongRoPE}

\label{sec:method}
Building on these observations, we propose LongRoPE. This approach initially employs a novel and efficient search algorithm that takes advantage of both types of non-uniformities, and subsequently utilizes this mechanism to extend the LLM context window to exceed 2 million tokens.

\subsection{Problem Formulation}

The presence of these two forms of non-uniformity significantly expands the solution space and adds layers of complexity to the optimization process. To manage this, we conceptualize the multidimensional non-uniform position interpolation optimization task as a search problem.

Given an LLM designed for a context window of length $L'$ and input documents $\mathbf{X}$ where each $\mathbf{x}\in\mathbf{X}$ contains more tokens than $L'$, we represent the original rotary angle used in the RoPE embedding for the $i^{th}$ dimension at position $n$ as $\frac{n}{\beta^i}$. The corresponding optimization objective can be stated as:

\begin{equation}
\begin{gathered}
		\label{eq:problem}

\underset {\mathbf{x}\in\mathbf{X}; \, |\mathbf{x}|\ge L'}{ \operatorname {arg\,min} } \,  \mathcal{L}\, \left( \text{LLM} (\text{RoPE}, \mathbf{X})\right),	\text{where \;} 
\\
\scalebox{0.93}{
$\underset {\begin{subarray}{l} i=0,\cdot\cdot,\frac{d}{2}-1;\\ n\in[ 0,|\mathbf{x}|);\end{subarray}}{\text{RoPE}_(n)}= {\left[\cdots, \cos\left(\mathbb{I}(\hat{\lambda}_{i}, \hat{n})\times\frac{n}{\beta^i}\right),  \sin\left(\mathbb{I}(\hat{\lambda}_{i}, \hat{n})\times\frac{n}{\beta^i}\right),\cdots\right] }$}\\
	\text{where \;} \mathbb{I}(\hat{\lambda}_{i},\hat{n})=\begin{cases}
	1&\mbox{\;  $n<\hat{n}$}\\
{\frac{1}{\lambda}_{i}}&\mbox{\; $n\geq\hat{n}$}
\end{cases}
	\end{gathered}
\end{equation}

In this context, we define a collection of rescaling factors, $\mathbb{I}(\hat{\lambda}_{i},\hat{n})$, to address two types of non-uniformity. Here, $\hat{\lambda}_i$ characterizes the variation among RoPE dimensions, while $\hat{n}$ corresponds to the variation with respect to token positions. The function $\mathbb{I}(\hat{\lambda}_{i},\hat{n})$ is utilized to adjust the rotation angle for the $i^{th}$ RoPE dimension. Specifically, $\hat\lambda_{i}$ serves as the scaling parameter, and $\hat{n}$ is the threshold for token positions. For token positions before $\hat{n}$ (i.e., $n < \hat{n}$), the rescaling is not applied, and the standard RoPE rotation angle, $\frac{n}{\beta^i}$, is maintained. Once the token position reaches or exceeds $n \geq \hat{n}$, the rescale factor comes into effect.

Assuming a desired context window length of $L'$, our aim is to determine the set of optimal rescaling factors ($\mathbb{I}(\hat{\lambda}_{0},\hat{n})$, $\mathbb{I}(\hat{\lambda}_{1},\hat{n})$, ..., $\mathbb{I}(\hat{\lambda}_{i},\hat{n})$, ...) across the $1^{st}$ through $d^{th}$ dimensions of the RoPE. Consequently, by applying these rescaled $\text{RoPE}$ factors, the target $\text{LLM}$ can attain the lowest possible value of next token prediction loss, $\mathcal{L}$ (corresponding to perplexity), when evaluated on input sequences $\mathbf{X}$ having token length $L'$.

\subsection{Searching the Non-uniform Position Interpolation}

To address the issue presented in Eq.~\ref{eq:problem}, we propose a straightforward yet powerful approach that optimizes RoPE rescale factors, thereby effectively leveraging the positional embedding’s multidimensional and non-uniform properties.

\textbf{Search space}. Our method defines an extensive search space that encompasses both types of non-uniformities, as detailed in Table~\ref{tbl:searchspace}. Concretely, we enable the determination of individual rescale factors for each dimension within the RoPE embedding. For greater simplicity in search space configuration, the method searches over $\lambda_i$ and $\hat{n}$, rather than directly exploring $\mathbb{I}(\hat{\lambda}_{i},\hat{n})$, where we have $\hat{\lambda}_{i}=1/\lambda_i$. Table~\ref{tbl:searchspace} demonstrates that $\lambda_i$ is permitted to vary from a baseline of 1.0 (corresponding to direct extrapolation) up to an upper bound of $s\times1.25$ (providing more pronounced interpolation than PI), utilizing an increment of 0.01 at each step, with $s$ representing the intended context window extension factor.

The parameter $\hat{n}$ determines how many of the earliest token positions preserve the original RoPE embedding without applying position interpolation. In practice, we let $\hat{n}$ vary within the set \{0, 1, 2, 4, 8, 12, 16, 20, 24, 28, 32, 64, 128, 256\}. Setting $\hat{n}=0$ means that rescale factors obtained from the search are applied across all token positions.

\textbf{Evolution-based search}. The search space detailed in Table~\ref{tbl:searchspace} encompasses a vast array of positional interpolation possibilities, making thorough exploration highly demanding. For instance, achieving a $s=4\times$ extension results in $400^{128/2}\times14 = 4\times10^{167}$ potential configurations. As the extension ratio increases, the search space grows at an exponential rate. To manage this complexity, we adopt evolution search~\cite{guo2020single} and implement two additional optimization strategies that significantly enhance search efficiency. The complete search workflow is outlined in Algorithm~\ref{alg:search}.

\textit{Optimized initial population generation}. Rather than relying solely on random initialization for the population of $P$ rescale factors, we explicitly insert the three RoPE rescale factors associated with PI, NTK, and YaRN as fixed individuals within the starting population. To fill the remaining $P$-3 slots, we generate individuals by applying random mutations to these three baseline rescale factors, each mutation occurring with a probability $p$.

\textit{Monotonically non-decreasing constraint}. Once the initial population is produced, we evaluate each individual by calculating its perplexity under the LLM. This involves applying the respective RoPE rescale factors within the LLM to determine the perplexity for the input $\mathbf{X}$. The $k$ individuals with the lowest perplexity values are selected as parents for subsequent evolutionary steps. Owing to the extremely large search space, straightforward mutation and crossover operations risk producing many low-quality candidates, which can result in a significant number of redundant perplexity evaluations. This inefficiency is exacerbated for large $L'$, where each perplexity calculation demands substantial computational resources.

To address this issue, we enforce a monotonic non-decreasing constraint on the set of sampled RoPE rescaling factors, specifically $\lambda_i \le \lambda_{i+1}$. Only those RoPEs that meet this requirement are utilized in LLM perplexity assessments, which greatly narrows the search space. In particular, our approach stipulates that $\lambda_i$ should strictly increase with respect to the RoPE dimension ($i=0,\ldots,63$). This requirement draws upon insights from NTK theory~\cite{jacot2018neural,tancik2020fourier,ntk}, which indicates that lower-indexed dimensions (with higher frequencies) necessitate less interpolation (hence smaller $\lambda_i$), while higher-indexed, lower-frequency dimensions can tolerate greater interpolation (thus allowing larger $\lambda_i$).

\begin{algorithm}[t]
	\small
	\caption{The search algorithm}
	\textbf{Input:} target LLM, input samples $\mathbf{X}$, population size $P$, mutation size $N_1$, crossover size $N_2$, max iterations $\mathcal{T}$, mutate probability $p$\\

	\begin{algorithmic}[1]
		\label{alg:search}
		\STATE Initialize Top-k to be empty;
		\STATE Set $\text{P}_0$ as the output of \textit{Initialize\_population\_with\_optimization} ($P$, $\mathbf{X}$, $p$);
		\FOR{$i$=$1$ to $\mathcal{T}$}
		\STATE Run \textit{Compute\_perplexity} (LLM, $\text{P}_{i-1}$, $\mathbf{X}$);
		\STATE Update Top-k via \textit{Update\_Topk} (Top-k, $\text{P}_{i-1}$);
		\STATE Generate $\text{P}_{mutation}$ using \textit{Mutation\_with\_mono\_constraint} (Top-k, $N_1$, $p$);
		\STATE Create $\text{P}_{crossover}$ through \textit{Crossover\_with\_mono\_constraint} (Top-k, $N_2$);
		\STATE Construct $\text{P}_i$ by taking the union of $\text{P}_{mutation}$, $\text{P}_{crossover}$, and Top-k;
		\ENDFOR
		\STATE Output the member of Top-k with the minimum perplexity value;
	\end{algorithmic}
\end{algorithm}

\textbf{8$\times$ extension without fine-tuning}. Through evolutionary search, our approach efficiently discovers optimal non-uniform RoPE rescale factors, selectively maintaining critical dimensions and positions in order to reduce information degradation resulting from interpolation. As shown in Fig.\ref{fig:non-ft}, our technique successfully expands LLaMA2's context window from 4k to 32k tokens, all without requiring fine-tuning. Conversely, other strategies like PI, as well as non-uniform NTK and YaRN, exhibit a marked increase in perplexity beyond a 2$\times$ extension.

\subsection{Extending LLM Context Window to 2048K}
\label{sec:extendto2048k}

\textbf{Progressive extension to 2048k}. In this section, we present an approach for expanding the context window of pre-trained LLMs from the standard 4k tokens up to more than 2048k tokens. As previously shown, applying our non-uniform positional interpolation method enables an 8$\times$ increase without requiring any fine-tuning. However, for substantially larger expansions such as 512$\times$, fine-tuning becomes indispensable. A potential strategy involves identifying suitable RoPE rescaling factors corresponding to the target context window of 2048k and performing subsequent fine-tuning. Nonetheless, this procedure is hindered by the significant computational resources required for training. Additionally, our observations indicate that effective fine-tuning at such large extension ratios is difficult to achieve (refer to the Appendix for details).

Fortunately, LongRoPE demonstrates efficacy with both pre-trained and further fine-tuned extended LLMs. Consequently, we propose an incremental and efficient approach capable of reaching the target length of 2048k tokens using only 1k steps of fine-tuning, all conducted within a maximum training length of 256k.

$\Diamond$ \textit{LongRoPE search applied to pre-trained LLMs for 256k extension}. Using LLaMA2 as a case study, we perform a search with desired context window lengths of 128k and 256k. At this point, the scaling factors are 32$\times$ and 64$\times$, accordingly.

$\Diamond$ \textit{Fine-tuning to 256k}. Next, we adapt the pre-trained LLM for a context window length of 256k through fine-tuning. To do this, we initially train LLaMA2 for 400 steps utilizing RoPE rescaled factors set for 128k. Afterwards, we swap the RoPE rescaled factors to those appropriate for 256k on the resulting checkpoint and continue training for an extra 600 steps. This staged approach demonstrates greater efficiency compared to performing the fine-tuning to 256k from the outset.

$\Diamond$ \textit{Scaling fine-tuned extended LLM to 2048k using LongRoPE search}. In the concluding stage, we conduct an additional search procedure on the already fine-tuned LLM set at 256k context length. This enables the model to support a massive context window of 2048k, all without the need for further fine-tuning. The total achieved extension factor is $512\times$.

\textbf{Recovery within reduced context window}. When scaling the context window up to an extensive 2048k tokens, we observe that performance deteriorates in the initial, standard context window range. This degradation aligns with established observations in positional interpolation~\cite{pi}, since mapping positions to higher dimensions within the unchanged context window compresses them into a much smaller subspace, which impairs language model effectiveness. Notably, applying a 512$\times$ extension ratio leads to significant congestion of positions in the canonical 4k context window.

To address this issue, we conduct an additional evolutionary search on the expanded LLM to fine-tune the RoPE rescale factors tailored to shorter context lengths (such as 4k and 8k). For these shorter sequences, the upper limit for searched $\lambda$ is decreased, as there is a reduced need for positional interpolation. At inference time, the LLM automatically modifies the relevant RoPE rescale factors in accordance with the current context length.

\section{LongRoPE}

\label{sec:method}
Building on these observations, we propose LongRoPE, which initially employs an optimized search strategy to take full advantage of both forms of non-uniformity, and subsequently leverages this to expand the LLM context window to more than 2 million tokens.

\subsection{Problem Formulation}

The aforementioned non-uniformities contribute to an expansive solution landscape and complicate the optimization procedure. To tackle this challenge, we recast the multidimensional non-uniform position interpolation optimization as a search-based problem.

Consider an LLM with a designated context window of size $L'$ and a collection of lengthy input texts $\mathbf{X}$, where each sequence $\mathbf{x} \in \mathbf{X}$ contains more tokens than $L'$. Let the initial rotary angle for the $i^{th}$ dimension in the RoPE embedding at token position $n$ be defined as $\frac{n}{\beta^i}$. The corresponding optimization objective is described below:

\begin{equation}
\begin{gathered}
		\label{eq:problem}

\underset {\mathbf{x}\in\mathbf{X}; \, |\mathbf{x}|\ge L'}{ \operatorname {arg\,min} } \,  \mathcal{L}\, \left( \text{LLM} (\text{RoPE}, \mathbf{X})\right), \ \text{where} 
\\
\scalebox{0.93}{
$\underset {\begin{subarray}{l} i=0,\ldots,\frac{d}{2}-1;\\ n\in[ 0,|\mathbf{x}|);\end{subarray}}{\text{RoPE}_(n)}= {\left[\ldots, \cos\left(\mathbb{I}(\hat{\lambda}_{i}, \hat{n}) \frac{n}{\beta^i}\right), \sin\left(\mathbb{I}(\hat{\lambda}_{i}, \hat{n}) \frac{n}{\beta^i}\right),\ldots \right]}$}\\
\text{where} \; \mathbb{I}(\hat{\lambda}_{i},\hat{n}) = 
\begin{cases}
    1 &\text{if}\quad n<\hat{n}\\
    \frac{1}{\lambda_{i}} &\text{if}\quad n\geq\hat{n}
\end{cases}
\end{gathered}
\end{equation}

In this approach, we define a collection of rescale factors, $\mathbb{I}(\hat{\lambda}_{i},\hat{n})$, to address both types of non-uniformity. Here, $\hat{\lambda}_i$ refers to the adjustment applied to the $i^{th}$ RoPE dimension, and $\hat{n}$ represents the threshold token position related to non-uniformity. The function $\mathbb{I}(\hat{\lambda}_{i},\hat{n})$ is utilized to modify the rotation angle corresponding to the $i^{th}$ RoPE dimension; $\hat{\lambda}_i$ serves as the scaling parameter, while $\hat{n}$ acts as a cutoff for token positions. For the first $\hat{n} - 1$ token positions, the rescale factor $\hat{\lambda}_i$ is not employed, resulting in use of the standard RoPE rotary angle $\frac{n}{\beta^i}$. Once token positions reach $n\ge\hat{n}$, the scaling factor is then incorporated.

With a desired context window length of $L'$, the task is to determine the most effective set of rescale coefficients ($\mathbb{I}(\hat{\lambda}_{0},\hat{n})$, $\mathbb{I}(\hat{\lambda}_{1},\hat{n})$, ...$\mathbb{I}(\hat{\lambda}_{i},\hat{n})$...) across each RoPE dimension from $1$ to $d$. Consequently, applying these rescaled $\text{RoPE}$ parameters should enable the target $\text{LLM}$ to achieve the lowest possible next-token prediction loss $\mathcal{L}$ (equivalent to minimizing perplexity) for sample inputs $\mathbf{X}$ of length $L'$.

\subsection{Searching the Non-uniform Position Interpolation}

To address the challenge posed in Eq.~\ref{eq:problem}, we propose a straightforward yet robust approach that identifies optimal RoPE rescale factors to take full advantage of the heterogeneous characteristics present across different dimensions in position embeddings.

\textbf{Search space}. Our method establishes an extensive search space that explicitly incorporates both types of non-uniformity. The structure of this search space is detailed in Table~\ref{tbl:searchspace}. In particular, we enable individualized rescale factors to be determined for every dimension in the RoPE embeddings. For simplicity in configuring the search space, we conduct the search over $\lambda_i$ and $\hat{n}$ rather than directly optimizing $\mathbb{I}(\hat{\lambda}_{i},\hat{n})$, where $\hat{\lambda}_{i}=1/\lambda_i$. According to Table~\ref{tbl:searchspace}, each $\lambda_i$ may be selected from a minimum value of 1.0 (representing basic extrapolation) to a maximum of $s\times1.25$ (corresponding to interpolation regions even beyond those used by PI), incremented in steps of 0.01. Here, $s$ refers to the context window extension ratio being targeted.

The parameter $\hat{n}$ determines how many of the earliest token positions keep their original RoPE embedding, foregoing position interpolation. In practice, we permit $\hat{n}$ to vary over the set \{0, 1, 2, 4, 8, 12, 16, 20, 24, 28, 32, 64, 128, 256\}. If $\hat{n}=0$, the rescale factors found by the search are applied to every token position.

\textbf{Evolution-based search}. The extensive search space outlined in Table~\ref{tbl:searchspace} encompasses a multitude of positional interpolation configurations, making exhaustive exploration highly impractical. For instance, with an extension of $s=4\times$, the possible choices explode to $400^{128/2}\times14=4\times10^{167}$. As the extension ratio grows, this complexity increases exponentially. To efficiently navigate this vast landscape, we employ evolution search~\cite{guo2020single} alongside two optimization strategies designed to significantly enhance the speed of the search process. The overall procedure is depicted in Algorithm~\ref{alg:search}.

\textit{Optimized initial population generation}. In place of generating all $P$ rescale factors at random, we directly incorporate the three RoPE rescale factors related to PI, NTK, and YaRN as explicit members of the initial population. For the other $P$-3 individuals, these three rescale factors are subjected to random mutations, each with a probability of $p$.

\textit{Monotonically non-decreasing constraint}. Once the initial population is established, we evaluate the LLM perplexity for every individual. To do so, we implement the designated RoPE rescale factors within the target LLM and determine the perplexity for the input $\mathbf{X}$. Individuals with the lowest perplexity scores—specifically the top-k—are selected as parents for the subsequent evolutionary process. Due to the enormous size of the search space, straightforward mutation and crossover operations may frequently lead to suboptimal candidates, resulting in redundant perplexity evaluations. This inefficiency becomes especially pronounced when $L'$ is large, since each perplexity evaluation requires a significant amount of inference time.

To mitigate this issue, we enforce a non-decreasing monotonicity condition on the rescaled RoPE factors sampled, such that $\lambda_i \le \lambda_{i+1}$. RoPE configurations adhering to this requirement are the only ones selected for perplexity assessments in the LLM, which serves to notably reduce the computational burden associated with the search process. In particular, we stipulate that $\lambda_i$ exhibits a monotonic increase across RoPE dimensions (that is, for $i = 0,\ldots,63$). This constraint is motivated by insights from NTK theory~\cite{jacot2018neural,tancik2020fourier,ntk}, which implies that lower-dimensional components associated with higher frequency content necessitate less interpolation (corresponding to smaller $\lambda_i$ values), while higher-dimensional components with lower frequency can tolerate more interpolation (resulting in larger $\lambda_i$ values).

\begin{algorithm}[t]
	\small
	\caption{The search algorithm}
	\textbf{Input:} target LLM, input samples $\mathbf{X}$, population size $P$, mutation size $N_1$, crossover size $N_2$, maximum number of iterations $\mathcal{T}$, mutation probability $p$\\

	\begin{algorithmic}[1]
		\label{alg:search}
		\STATE Top-k $\leftarrow \phi$;	
		\STATE $\text{P}_0$ $\leftarrow$ \textit{Initialize\_population\_with\_optimization} ($P$, $\mathbf{X}$, $p$);
		\FOR{$i = 1$ to $\mathcal{T}$}
		\STATE \textit{Compute\_perplexity} (LLM, $\text{P}_{i-1}$, $\mathbf{X}$);
		\STATE Top-k $\leftarrow$ \textit{Update\_Topk}(Top-k, $\text{P}_{i-1}$);
		\STATE $\text{P}_{mutation}$ $\leftarrow$ \textit{Mutation\_with\_mono\_constraint}(Top-k, $N_1$, $p$); 
		\STATE $\text{P}_{crossover}$ $\leftarrow$ \textit{Crossover\_with\_mono\_constraint}(Top-k, $N_2$);
		\STATE $\text{P}_i$ $\leftarrow$ $\text{P}_{mutation}$ $\cup$ $\text{P}_{crossover}$ $\cup$ Top-k;
		\ENDFOR
		\STATE Output the member of Top-k exhibiting the minimum perplexity;
	\end{algorithmic}
\end{algorithm}

\textbf{8$\times$ extension without fine-tuning}. Through evolutionary search, we efficiently discover non-uniform RoPE rescale factors that prioritize maintaining essential dimensions and positions, thereby reducing the information loss caused by interpolation. As shown in Fig.\ref{fig:non-ft}, this strategy enables us to expand LLaMA2's context window from 4k to 32k tokens without the need for any fine-tuning. In comparison, previously established approaches like PI, as well as non-uniform NTK and YaRN, result in significant increases in perplexity after merely doubling the context window.

\subsection{Extending LLM Context Window to 2048K}
\label{sec:extendto2048k}

\textbf{Progressive extension to 2048k}. In this section, we present our approach for expanding the context window of pre-trained LLMs from the standard 4k up to more than 2048k. Our results indicate that utilizing non-uniform positional interpolation permits an 8$\times$ extension without the need for additional fine-tuning. However, scaling further, such as reaching 512$\times$ the original context length, necessitates fine-tuning. A common strategy involves identifying optimal RoPE rescale factors at the desired 2048k target before commencing fine-tuning. Yet, this process is often limited by the substantial computational resources it demands. Furthermore, our observations reveal that it is generally difficult to effectively fine-tune LLMs at such high extension rates (refer to Appendix).

Fortunately, LongRoPE demonstrates effectiveness on both base models and those further fine-tuned for extended context. Accordingly, we propose a streamlined, incremental approach that attains the desired 2048k context window using only 1k fine-tuning steps, all performed at a training length limited to 256k.

$\Diamond$ \textit{Scaling pre-trained LLMs to 256k with LongRoPE search}. Using LLaMA2 for illustration, we perform a search procedure for desired context window lengths of 128k and 256k. At this point, the expansion factors amount to 32$\times$ and 64$\times$, respectively.

$\Diamond$ \textit{Fine-tuning to 256k}. Subsequently, the pre-trained LLM is further optimized to reach a context window of 256k. In detail, we initially continue training LLaMA2 for 400 steps, employing RoPE rescaled factors corresponding to 128k. Afterwards, we substitute these rescaled factors with those for 256k on the resulting checkpoint and proceed with 600 additional fine-tuning steps. This staged approach is demonstrated to be more effective than directly performing fine-tuning at the 256k context window.

$\Diamond$ \textit{LongRoPE search enables 2048k extension of fine-tuned LLM}. As a final step, we apply a subsequent search process to the LLM that has already undergone fine-tuning at 256k sequence length. This approach successfully expands the model's context window up to 2048k tokens, eliminating the need for any additional fine-tuning. The overall expansion factor thus achieved is $512\times$.

\textbf{Restoration of performance in shorter context windows}. Upon increasing the context window to an extensive 2048k length, we observe diminished effectiveness within the model’s initial context window. This phenomenon is a recognized drawback of positional interpolation~\cite{pi}, where embeddings in higher dimensions are compressed into a limited region inside the original window, thereby degrading the language model’s overall capabilities. Specifically, applying a 512$\times$ expansion compresses positions in the initial 4k context window considerably, resulting in significant overlap.

To address this issue, we conduct an additional evolutionary search on the expanded LLM, fine-tuning the RoPE rescale parameters specifically for shorter context lengths (such as 4k and 8k). For these shorter ranges, we constrain the upper bound of the searched $\lambda$ since there is a decreased need for positional interpolation. At inference time, the LLM adaptively selects the appropriate RoPE rescale factors as needed.

 \section{Experiments}

\subsection{Setup}
\textbf{Evaluation Tasks and models}. LongRoPE is implemented on both LLaMA2-7B and Mistral-7B, with performance assessed along three key dimensions: (1) perplexity for long-context LLMs when processing lengthy documents; (2) performance on the Passkey retrieval task, which tests the model’s effectiveness at extracting a straightforward passkey hidden among predominantly unrelated text; and (3) results on established LLM benchmarks constrained to a context window of 4096 tokens.

\textbf{Fine-tuning}. For LLaMA2, our approach involves setting the learning rate at 2e-5, employing a linear decay schedule, and utilizing a global batch size of 32. Fine-tuning is performed over 400 steps with the Redpajama~\cite{redpajama} dataset, where the data is organized into 128k-token chunks that are wrapped with BOS and EOS tokens. Building on the checkpoint obtained from this phase, we further fine-tune for an extra 600 steps to extend to a 256k context window. Training for the 128k context window utilizes 8 A100 GPUs and leverages the distributed system described in~\cite{cube}, whereas the 256k context window necessitates 16 A100 GPUs. For Mistral, we adopt a fixed learning rate of 1e-6 together with a global batch size of 64. Both the 128k and 256k context window models are fine-tuned following the protocol established in YaRN~\cite{yarn}: 400 training steps on the Long-Data Collections provided by Together Computer~\cite{mistraldata} with sequences of length 16k. This process is executed on 4 A100 GPUs.

\textbf{Search}. When the target window size is up to 256k, we set $P$=64, $N_1$=$N_2$=16, $p$=0.3, and $\mathcal{T}$=40, and at each iteration we retain the top-32 candidates for mutation and crossover operations. Perplexity evaluation employs 5 randomly chosen samples from the PG19 validation set, each meeting the minimum length set by the desired context window. For target window sizes exceeding 512k, we reduce the population as well as the numbers used for mutation and crossover by half. In this regime, perplexity is computed on 3 randomly drawn samples from the Pile-Books3~\cite{pile} validation dataset.

\textbf{Baselines}. To achieve a context length of 2048k, we conducted fine-tuning of models on context windows of 128k and 256k tokens. As a result, we refer to the versions as LongRoPE-2048k (ft=128k) for LLaMA2 and LongRoPE-2048k (ft=256k) for Mistral. Our evaluation involves these four variants alongside leading context window extension methods, focusing particularly on publicly available LLMs that have been fine-tuned post-positional interpolation using techniques such as PI, NTK, and YaRN. The comparison set includes Together-32k~\cite{together}, Code LLaMA~\cite{codellama}, LongLoRA-full-FT-100k~\cite{longlora}, as well as YaRN-LLaMA and YaRN-Mistral~\cite{yarn}.

\subsection{Main Results}

\label{sec:mainresult}
\textbf{Language modeling for long sequences up to 256k}. Our analysis starts by benchmarking our approach against current leading extended LLMs on sequence lengths up to 256k. To illustrate the broad applicability of our methods, we select two datasets: the test split of Proof-pile~\cite{pg19} and PG19~\cite{pile}. Perplexity is assessed across multiple context sizes using a sliding window of 256 tokens. For the PG19 dataset, we utilize all 100 documents included in the test partition. In the case of Proof-pile, following the procedure outlined in YaRN~\cite{yarn}, we randomly sample 10 items, each with a minimum length of 128k.

Table~\ref{tbl:proof-pile} and Table~\ref{tbl:pg19} present a perplexity analysis of LLaMA2 and Mistral models, each enhanced using various interpolation strategies, across the Proof-pile and PG19 datasets, respectively. Two principal findings emerge: \textbf{(1)} our interpolated models consistently achieve lower perplexity scores as the evaluation context grows from 4k up to 256k tokens, demonstrating their enhanced capacity to utilize extended contexts effectively. \textbf{(2)} Remarkably, the LongRoPE-2048k variants maintain superior performance over leading baselines within 256k-token contexts, despite operating with a context window up to 16$\times$ larger—a scenario where short-context performance is often difficult to preserve.

\textbf{Long sequence language modeling beyond 2000k}. To assess performance on ultra-long documents, we employ the Books3~\cite{pile} dataset. For efficiency in evaluation, 20 books are randomly chosen, each with a length greater than 2048k, and we utilize a sliding window approach of 256k.

As indicated in Table~\ref{tbl:books3}, LongRoPE is able to extend the context window of both LLaMA2-7B and Mistral-7B models to 2048k tokens, while matching or surpassing baseline perplexity scores within shorter ranges of 8k to 128k. Distinct differences in performance are evident between the two models at the 2048k window: while Mistral outperforms all baselines for shorter input lengths, its perplexity rises above 7 for contexts longer than 256k. LLaMA2's perplexity follows the anticipated trend, generally decreasing as the context lengthens, with only slight increases seen at 1024k and 2048k. Additionally, for LLaMA2, fine-tuning LongRoPE-2048k at 256k yields superior results compared to fine-tuning at 128k; this is attributed to the smaller secondary extension ratio involved (8$\times$ versus 16$\times$). Conversely, Mistral achieves better outcomes when fine-tuned with a 128k window. This difference primarily stems from the use of a 16k training length for both 128k and 256k fine-tuning, following the YaRN protocol, which restricts Mistral's capability to further increase its context window after the fine-tuning phase.

\textbf{Passkey retrieval}. Next, we investigate how the usable context window influences generation tasks. Adopting the synthetic passkey retrieval benchmark from~\cite{passkey}, we task the model with locating a randomly embedded five-digit number (“passkey”) within an extended sequence of text. The structure of the prompt employed is provided in the appendix. For the evaluation, we conduct 10 trials of passkey retrieval, each time positioning the passkey uniformly at random within the entire length of the evaluation context.

Fig.~\ref{fig:passkey} presents a comparison of retrieval accuracy between our approach and baseline models. The accuracy of previous models plummets to zero for context lengths beyond 128k tokens. Conversely, LongRoPE-LLaMA2-2048k (ft=256k) demonstrates strong performance on the challenging task of passkey retrieval among millions of tokens, consistently achieving a high retrieval accuracy ($\ge$90\%) from 4k up to 2048k. Meanwhile, LongRoPE-Mistral-2048k (ft=128k) maintains perfect (100\%) retrieval up to 1800k tokens, with accuracy decreasing to 60\% at 2048k—consistent with the observed increase in perplexity at 2048k reported in Table~\ref{tbl:books3}.

\textbf{Evaluation on standard benchmarks within the original context window}. LongRoPE-2048k models are assessed on their performance within the initial context window, employing the Hugging Face Open LLM Leaderboard~\cite{huggingfaceboard} in both zero-shot and few-shot paradigms. The following benchmark datasets are utilized: ARC-Challenge~\cite{allenai:arc} with 25 shots, HellaSwag~\cite{zellers2019hellaswag} with 10 shots, MMLU~\cite{mmlu} with 5 shots, and TruthfulQA~\cite{lin2021truthfulqa} in a zero-shot configuration.

According to Table~\ref{tbl:commonsense}, our models perform on par with the original benchmark, which was developed for a shorter context window, and surpass the initial Mistral's performance on TruthfulQA by +0.5\%. While LongRoPE-LLaMA2-2048k—fine-tuned with a 256k context length—exhibits somewhat greater drops in performance, its results generally stay within acceptable limits across the majority of tasks.

\subsection{Ablation Results}

\textbf{Effectiveness of the second positional interpolation}. Within the framework of our progressive extension method, we leverage our search procedure to apply a secondary, non-uniform positional interpolation to the fine-tuned, extended LLMs. To assess this approach, we perform experimental evaluation using our fine-tuned LLaMA2-256k model, scaling it further to 512k, 1024k, and 2048k with both PI and YaRN. As indicated in Table~\ref{tbl:secondsearch}, our non-uniform positional interpolation maintains stable perplexity across different extension sizes. In comparison, the perplexity observed with PI and YaRN escalates rapidly as the extension ratio increases.

\textbf{Effectiveness of recovery at shorter context lengths.} In order to address reduced performance at shorter context windows, we recalibrate the RoPE factors in LongRoPE-2048k using our search methodology. In particular, we impose stricter limits on the maximum scale factors during the search, which reduces the degree of interpolation for short context lengths of 4k and 8k. Table~\ref{tbl:shortperformance} provides a perplexity comparison for LongRoPE-LLaMA2-2048k on the Proof-pile dataset at these lengths, alongside overall LLM benchmark accuracy. These results indicate a marked improvement in performance at shorter context lengths.

\textbf{Assessment of the Two Types of Non-uniformities.} To understand the specific contribution of each type of non-uniformity, we conduct ablation studies. Our approach involves two experiments: (i) we expand LLaMA2-7B to short context lengths of 16k and 32k via various strategies—namely PI, optimizing solely over the RoPE dimension, and jointly optimizing both non-uniformities; (ii) we further scale our fine-tuned 256k context LLaMA2 model up to 2048k, applying analogous methods. Evaluation is based on perplexity measured without further fine-tuning. Referring to Table~\ref{tbl:searchdimension}, optimizing the non-uniformity within the RoPE dimension yields a notable decrease in perplexity compared to the linear interpolation technique used in PI. Adjusting token position non-uniformity also notably enhances results at the 16k and 32k context lengths, though this benefit does not persist at the 2048k context, which may be attributed to the vast sequence length. Simply retaining the initial sequence tokens and omitting interpolation appears ineffective in this regime, a matter that we aim to explore further in future studies.

\section{Related Works}

In addition to methods based on position interpolation, this section discusses related works of other approaches. 

\textbf{Retrieval-based} methods employ an external memory component to store extended historical context, alongside retrieval mechanisms that locate pertinent documents during inference~\cite{longllama,wang2023augmenting,borgeaud2022improving}. Implementing these methods often requires explicit architectural changes to the underlying LLMs. Differing from these approaches, our method introduces only minimal adjustments to positional embeddings, making it comparatively lightweight. Additionally, our framework extends to various long-context scenarios beyond retrieval, including tasks such as long document summarization and few-shot learning.

\textbf{Attention-based context window extensions}. In addition to extending context windows via positional embedding interpolation, various studies have expanded the input context using the initial LLM context window size through modifications of attention mechanisms~\cite{han2023lminfinite, streamingllm,parallelwindow}. Central to these approaches is the use of innovative attention masks to address the challenge of attention explosion introduced by additional positions. These methods function in a way that complements positional interpolation techniques.

\textbf{Fine-tuning based approaches} investigate strategies for adapting pre-trained LLMs to accommodate extended contexts through modifications to position embeddings. Studies such as Code LLaMA~\cite{codellama}, LLaMA2 Long~\cite{xiong2023effective}, and ScaledRoPE~\cite{liu2023scaling} select a significantly increased base value for RoPE, subsequently fine-tuning on the intended target sequence length. In contrast, our approach supports a range of target lengths and is capable of handling inputs exceeding 2M tokens. With the recognition that fine-tuning for extensive context lengths (greater than 128k) is computationally intensive and requires considerable GPU capacity, recent methods like LongLoRA~\cite{longlora} and PoSE~\cite{zhu2023pose} have been introduced to alleviate such demands. Our approach can be used in conjunction with these resource-efficient fine-tuning techniques.

\section{Conclusion}

This paper introduces LongRoPE, an approach that pushes the context length of LLMs to a new high of 2048k, all while preserving performance within their initial, shorter context ranges. By leveraging two types of non-uniformity inherent in RoPE positional encoding and applying an efficient evolutionary search algorithm, this technique delivers two main advantages: it not only supplies an optimal starting point for subsequent fine-tuning, but also supports an eightfold context window expansion without necessitating any fine-tuning at all. Furthermore, the method employs a stepwise expansion framework, utilizing 256k-context fine-tuned LLMs to ultimately extend the context window to 2048k without additional fine-tuning stages. Comprehensive experimental results support the efficacy of LongRoPE. We anticipate that LongRoPE-2048k will unlock a range of novel applications for long-context processing and stimulate more research in this area.

\section*{Broader Impacts} The objective of this paper is to contribute to progress within the domain of Machine Learning. While the research may carry a range of possible implications for society, we do not identify any particular issues that warrant special mention at this time.

	\balance

\end{document}